% !TeX encoding = UTF-8
% !TeX program = lualatex
\documentclass[a4paper,14pt]{diplom}

%%%%%%%%%%%%%%%%%%%%%%%%%%
% Преамбула
%%%%%%%%%%%%%%%%%%%%%%%%%%

% --- Дополнительные пакеты ---
\usepackage{tikz}
\usepackage{listings}
\usepackage{booktabs}
\usepackage{algorithm2e}
\usepackage{icomma}
\usepackage{microtype}

% --- Настройки математики ---
\mathtoolsset{showonlyrefs}
\newcommand{\N}{\mathbb{N}}
\newcommand{\Z}{\mathbb{Z}}
\newcommand{\R}{\mathbb{R}}
\DeclareMathOperator{\sgn}{sgn}
\newcommand{\from}{\colon}
\renewcommand{\le}{\leqslant}
\renewcommand{\leq}{\leqslant}
\renewcommand{\ge}{\geqslant}
\renewcommand{\geq}{\geqslant}
\renewcommand{\emptyset}{\varnothing}
\renewcommand{\epsilon}{\varepsilon}

% --- Данные для титульного листа ---
\LetterHead{\bfseries Министерство науки и высшего образования Российской Федерации 
	Федеральное государственное бюджетное образовательное учреждение высшего образования <<Ярославский государственный университет имени П.\,Г.~Демидова>>}
\Kafedra{\bfseries Кафедра информационных и сетевых технологий}
\ZavKaf{Заведующий кафедрой,\\ к.\,ф.-м.\,н.,}{Д.\,Ю.~Чалый}
\DocumentType{{\fontsize{16pt}{18pt}\selectfont Выпускная квалификационная работа}}
\Title{{\fontsize{16pt}{18pt}\selectfont\bfseries Название работы}}
\Napr{{\fontsize{12pt}{14pt}\selectfont по направлению подготовки 09.04.03\\Прикладная информатика, корпоративные информационные системы}}
\Chief{Научный руководитель\\ к.\,ф.-м.\,н., доцент}{И.\,О.~Фамилия}
\Author{Студент~группы~ПИЭ\nobreakdash-23МО}{И.\,О.~Фамилия}

% Ключевые слова для реферата
\Keywords{информатика, прикладная математика, \LaTeX}

% --- Библиография ---
\addbibresource{references.bib}

%%%%%%%%%%%%%%%%%%%%%%%%%%
% Конец преамбулы
%%%%%%%%%%%%%%%%%%%%%%%%%%

\begin{document}

% Создаем титульный лист
\maketitle

% Реферат
\begin{abstract}
Это пример оформления дипломной работы с помощью издательской системы \LaTeXe.

Реферат размещается <<непосредственно>> за титульным листом. Объем реферата должен составлять не более половины страницы. В~реферате указываются параметры ВКР: объем работы в страницах, количество глав, иллюстраций, таблиц, приложений, использованных источников. Перечень ключевых слов должен включать от~5 до~15 слов или словосочетаний из текста работы, которые в~наибольшей мере характеризуют ее содержание и~обеспечивают возможность информационного поиска. Ключевые слова приводятся в именительном падеже и~печатаются прописными буквами полужирным шрифтом в~строку через запятые.

Текст реферата должен отражать объект исследования, цель работы, результаты работы, область применения, степень внедрения или рекомендации по~внедрению.
\end{abstract}

% Содержание
\tableofcontents[Содержание]

% Введение (ненумерованная глава)
\chapternonum{Введение}
Во введении обосновывается актуальность выбранной темы,
описываются объект и~предмет исследования, цели и~задачи, методы исследования и~приводится краткое описание структуры работы.

% Основная часть
\chapter[Структура выпускной квалификационной работы]{Структура выпускной\\ квалификационной работы}
Выпускная квалификационная работа включает следующие структурные элементы:
\begin{enumerate}[label=\arabic{enumi})]
  \item титульный лист;
  \item реферат;
  \item содержание;
  \item введение;
  \item основную часть:
  \begin{itemize}
    \item глава 1,
    \item глава 2,
    \item \dots;
  \end{itemize}
  \item заключение;
  \item список использованных источников (список литературы);
  \item приложения.
\end{enumerate}
Каждый структурный элемент ВКР начинается с новой страницы.

Разделы <<Введение>> и <<Заключение>> не нумеруются.
В~них не~должно содержаться рисунков, формул и~таблиц.

Основная часть выпускной квалификационной работы не~требует специального заголовка, а~делится на~главы, состоящие из~параграфов, которые в~свою очередь, могут быть разбиты на~пункты. Каждая из~этих составляющих имеет заголовок, входящий в~состав содержания. Слова <<глава>>, <<параграф>>, <<пункт>> в~заголовках не~используются.
Нумерация выше названных составляющих основной части производится по~числовой иерархической системе, причем после последней цифры, а~также после заголовка точка не~ставится.

\chapter{Оформление элементов текста}

\section{Рисунки и таблицы}
Иллюстрации (фотографии, рисунки, чертежи, графики, диаграммы и~т.\,п.) обозначаются сокращенно словом <<Рис.>>, которое пишется под иллюстрацией  с~прописной буквы и~выделяется полужирным шрифтом. Нумеруются иллюстрации арабскими цифрами. Нумерация сквозная по~всему тексту ВКР. Пример "--- рис.~\ref{fig:1}.
Под рисунком по~центру размещаются его наименование и~поясняющие надписи.
Иллюстрации располагают сразу~же после ссылки на них в~тексте ВКР.

% Рисунок
\begin{figure}[!ht]
  \centering
  \includegraphics[width=0.3\textwidth]{apple.jpg}
  \caption{Название рисунка}
  \label{fig:1}
\end{figure}

Таблицы нумеруются в рамках раздела арабскими цифрами. Слово <<Таблица>> и~ее~номер пишется вверху, с~правой стороны над таблицей.
Ниже слова <<Таблица>>  посередине строки помещают ее название. Название  таблицы должно отражать ее~содержание, быть точным и~кратким. Название таблицы записывается с~прописной буквы и~выделяется полужирным шрифтом. Заголовки строк и~столбцов выделяются полужирным шрифтом.
Пример "--- табл.~\ref{tab:1}.

% Таблица
\begin{table}[tbh]
  \caption{Название таблицы}
  \label{tab:1}
  \centering
  \begin{tabular}{lccc}
    \toprule
    & \multicolumn{3}{c}{\textbf{Число глав}} \\
    \cmidrule(l){2-4}
    \textbf{Тип работы} & \textbf{Одна} & \textbf{Две} & \textbf{Три} \\
    \midrule
    Курсовая & 3 & 2 & 1 \\
    Работа бакалавра & 2 & 4 & 3 \\
    Диплом & 1 & 5 & 6 \\
    Магистерская диссертация & 0 & 4 & 5 \\
    \bottomrule
  \end{tabular}
\end{table}

\section{Заголовки и приложения}

\subsection{Заголовки}
Заголовки должны четко и кратко отражать содержание соответствующих разделов, подразделов, пунктов.
Заголовок печатают, отделяя от~номера пробелом, с~заглавной буквы. Точка в~конце заголовка не~ставится.
Заголовки выделяют полужирным шрифтом.
В~заголовках следует избегать сокращений (за исключением общепризнанных аббревиатур). В~заголовке не~допускается перенос слова на~следующую строку и~подчеркивание слов.
Выравнивание заголовков выполняется по~левому краю или по~центру строки (единообразно во~всей работе) без абзацного отступа.
Расстояние между названием глав и~последующим текстом должно равняться двум межстрочным интервалам. Такое~же расстояние выдерживается между заголовками главы и~параграфа.

\subsection{Приложения}
В~виде приложений оформляется материал, дополняющий основную часть ВКР.
Приложения обозначают прописными буквами русского алфавита, начиная с~А, за исключением букв Ё, З, Й, О, Ч, Ь, Ы, Ъ.
Каждое приложение начинается с~новой страницы.  При~этом в~верхнем правом углу страницы приводят слово <<Приложение>>, записанное строчными буквами с~первой прописной, с~указанием номера приложения.
Название приложения располагается ниже его обозначения на~отдельной строке по~центру  строчными буквами с~первой прописной и~выделяется полужирным шрифтом.
Приложения должны иметь общую с~основной частью документа сквозную нумерацию страниц.
В~тексте ВКР должны быть даны ссылки на~все приложения.
Ссылки на~приложения в~тексте ВКР должны быть организованы в~строго нумерационном порядке.
Пример оформления "--- приложение~\ref{app:source}.

\section{Нумерация страниц}
Все страницы текста ВКР, включая его иллюстрации и~приложения, должны иметь сквозную нумерацию.
Титульный лист считается страницей №~1, но~номер на~нем не~проставляется.
Номера страниц проставляются арабскими цифрами внизу страницы в~ее~правом углу или по~центру. В~случае необходимости номер на~некоторых страницах может быть проставлен вручную.

\chapter{Формулы}
Издательская система \LaTeX~\cite{tex_yarsu} предлагает широкий спектр средств для набора математических формул.
Подробно эта тема освещается в многочисленных книгах (см., например,~\cite{oetiker_lshort, kotelnikov2004}).
Ниже приводится лишь несколько простых примеров.

Если параметр~$a$ равен нулю, а~$b \ne 0$, то~уравнение \(a x = b\) не~имеет корней.

Определим функцию $\sgn \from \R \to \N$ следующим образом:
\begin{equation}
\sgn(x) = \begin{cases}
    1, & \text{при } x \ge 0,\\
   -1, & \text{иначе.}
\end{cases}
\label{eq:sgn}
\end{equation}
Из формулы~\eqref{eq:sgn} следует, что $x \sgn(x) = |x|$.

\emph{Множество рациональных чисел:}
\[
  \mathbb{Q} \coloneqq \Set*{ \frac{n}{m} \given n\in\Z, \ m\in\N}.
\]

Вероятность события $A$ при условии, что событие~$B$ произошло:
$\Pb{A \given B}$.

Вот так выглядит $(n\times k)$"~матрица:
\begin{equation}
M = \begin{pmatrix}
  m_{11} & m_{12} & \dots & m_{1k} \\
  m_{21} & m_{22} & \dots & m_{2k} \\
  \vdots & \vdots & \ddots & \vdots \\
  m_{n1} & m_{n2} & \dots & m_{nk} \\
  \end{pmatrix}
  \label{eq:matrix}
\end{equation}

\chapter{Рисуем с помощью TikZ}
Пакет Ti\textit{k}Z предлагает удобные инструменты для рисования диаграмм, блок"=схем, графов, графиков функций и~т.\,п~\cite{tikz, kirutenko2014}.
При этом рисунки сохраняются в~векторной графике, а~для  надписей используется тот~же шрифт, что и~в~основном тексте.
Простейшие примеры использования этого пакета изображены на~рис.~\ref{fig:block} и~\ref{fig:sin}.

\usetikzlibrary{positioning}
\usetikzlibrary{arrows.spaced}
\usetikzlibrary{calc}
\usetikzlibrary{shapes.geometric}

\begin{figure}[ht]
  \begin{minipage}[b]{0.40\textwidth}
    \centering
    \begin{tikzpicture}[rect/.style={rectangle, thick, draw=black}, node distance=15mm and 10mm]
      \node (source) [draw=black, shape=diamond, shape aspect = 2, align=center] {Источник\\ данных};
      \node (result) [rect, rounded corners, right = of source] {Результат};
      \draw[->, >=stealth'] (source.north) to[bend left] (result);
      \node (C1) at ($(source.east)!.5!(result.west)$) {};
      \node (user) [draw=black, ellipse, below = of C1] {Наблюдатель};
      \draw[->, >=stealth'] (user.west) -| (source.south);
      \draw[<->, >=stealth'] (user) to (result);
    \end{tikzpicture}
    \caption{Блок-схема}\label{fig:block}
  \end{minipage}
  \hfill
  \begin{minipage}[b]{0.54\textwidth}
    \centering
    \begin{tikzpicture}[scale=1.0, >=stealth']
      \draw[->] (0, -1.3) -- (0, 1.5) node[right] {y};
      \draw[->] (-3.4, 0) -- (3.8, 0) node[above left] {x};
      \draw[thick, smooth, samples at = {-3.4,-3.2,...,3.8}] plot ({\x},{sin(\x*180/pi)});
      \draw[dashed] (0, 1.0) node[left] {1} -- (pi/2, 1.0) -- (pi/2, -0.05) node[below] {$\frac{\pi}{2}$};
      \node[below=0.1] at (pi, 0) {$\pi$};
      \node[below right] at (0, 0) {$O$};
    \end{tikzpicture}
    \caption{График  функции $y = \sin(x)$}
    \label{fig:sin}
  \end{minipage}
\end{figure}

\chapter{Псевдокод}
Пакет \href{https://www.ctan.org/pkg/algorithm2e}{algorithm2e}
предлагает широкий спектр инструментов для создания и оформления псевдокода.
Также имеется возможность делать ссылки на строки кода.
Например, в строке~\ref{alg:lDoWhile} алгоритма~\ref{alg:quicksort}
целиком содержится цикл типа \texttt{do-while}.

\SetAlgorithmName{Алгоритм}{Список алгоритмов}{}
\SetAlgoCaptionSeparator{.}
\DontPrintSemicolon
\SetKwProg{Proc}{Procedure}{}{end}
\SetKwProg{Fn}{Function}{}{end}
\SetKwInOut{Input}{Вход}
\SetKwInOut{Output}{Выход}
\SetKwRepeat{DoWhile}{do}{while}
\SetKwBlock{Loop}{loop}{endloop}

\begin{algorithm}
  \caption{Быстрая сортировка}
  \label{alg:quicksort}
  \SetKwArray{Array}{A}
  \SetKwData{Low}{b}
  \SetKwData{High}{e}
  \SetKwFunction{QuickSort}{QuickSort}
  \SetKwFunction{Partition}{Partition}
  \Input{ массив \Array, индексы начала \Low и конца \High сортируемого фрагмента}
  \Output{ массив \Array, отсортированный по возрастанию}
  \BlankLine
  \Proc{\QuickSort{\Array, \Low, \High}}{
    \If{$\Low < \High$}{
      $m \leftarrow$ \Partition{\Array, \Low, \High}\;
      \QuickSort{\Array, \Low, $m$}\;
      \QuickSort{\Array, $m+1$, \High}\;
    }
  }
  \BlankLine
  \Fn{\Partition{\Array, \Low, \High}}{
    $v \leftarrow$ \Array{\Low}\;
    $i \gets \Low - 1$\;
    $j \gets \High + 1$\;
    \Loop{
      \DoWhile{$\Array{$i$} < v$}{
        $i \gets i + 1$\;
      }
      \lDoWhile{$\Array{$j$} > v$}{$j \gets j - 1$} \label{alg:lDoWhile}
      \If{$i \ge j$}{\Return{$j$}\;}
      поменять местами \Array{$i$} и \Array{$j$}\;
    }
  }
\end{algorithm}

% Заключение (ненумерованная глава)
\chapternonum{Заключение}
В заключении подводятся итоги выполненной работы, рассказывается о~том, что удалось и~что не~удалось сделать, описываются перспективы продолжения исследований.

% --- Библиография ---
\printbibliography[title={Список литературы}]

% --- Приложения ---
\appendix

% Настройка листингов (общая для всех приложений)
\lstset{
  tabsize=4,
  formfeed=\newpage,
  extendedchars=true,
  basicstyle=\ttfamily,
  commentstyle=\rmfamily\itshape,
  stringstyle=\slshape,
  numbers=left,
  numbersep=1em,
  stepnumber=1,
  numberstyle=\footnotesize\color{black},
}

\chapter{Исходный код программы на C++}
\label{app:source}
\lstdefinestyle{cpp}{
  language=[ANSI]C++,
  morekeywords={string, list}
}
% Загружаем код из файла sort.cpp
\lstinputlisting[style=cpp]{sort.cpp}

% Второе приложение – включаем из внешнего файла
\input{appendix2}

\end{document}

%%%%%%%%%%%%%%%%%%%%%%%%%%
% Конец документа
%%%%%%%%%%%%%%%%%%%%%%%%%%
